%==============================================================================
%  PLAN DE LA OBRA
%==============================================================================
\chapter*{Plan de la obra}
\addcontentsline{toc}{chapter}{Plan de la obra}
\markboth{PLAN DE LA OBRA}{PLAN DE LA OBRA}

La presente obra sistematiza de forma orgánica las instituciones de la \LIR{} y la
profunda reforma introducida por la \Lreforma{}, organizándose en trece capítulos que
combinan la dogmática de la insolvencia con el derecho procesal práctico aplicable ante
los tribunales ordinarios de justicia y ante la \Superir{}.

La estructura responde a un criterio deliberado: los capítulos \textsc{i} a \textsc{iii}
fijan las bases históricas, subjetivas y orgánicas del sistema; los capítulos \textsc{iv}
a \textsc{viii} recorren, en orden de intensidad decreciente de conservación de la
empresa, los cinco procedimientos concursales vigentes; los capítulos \textsc{ix} a
\textsc{xi} abordan los mecanismos de tutela de la masa y de control de la probidad del
deudor; y los capítulos \textsc{xii} y \textsc{xiii} proyectan el concurso hacia sus
fronteras laboral, penal e internacional.

Cada capítulo observa una arquitectura interna uniforme de cuatro bloques:
(i)~orientación teórica y justificación de la extensión; (ii)~estructura temática y
articulado de referencia; (iii)~manual práctico de tramitación paso a paso, con
flujogramas y modelos de escritos; y (iv)~doctrina y jurisprudencia esencial.

{\small
\begin{longtable}{@{}p{0.055\textwidth}p{0.34\textwidth}p{0.24\textwidth}p{0.30\textwidth}@{}}
\toprule
\textbf{Cap.} & \textbf{Título} & \textbf{Enfoque principal} & \textbf{Tratamiento} \\
\midrule
\endfirsthead
\toprule
\textbf{Cap.} & \textbf{Título} & \textbf{Enfoque principal} & \textbf{Tratamiento} \\
\midrule
\endhead
\bottomrule
\endfoot

\textsc{i} & Evolución histórica y dogmática del derecho concursal chileno
& De la punición a la reorganización
& Transición sistémica, principios rectores y crisis de la noción formal de insolvencia \\
\addlinespace
\textsc{ii} & El nuevo estatuto orgánico de los deudores tras la Ley \Nro 21.563
& Redefinición subjetiva del concurso
& Categorías de Empresa y Persona Deudora, y umbrales de tamaño de la Ley \Nro 20.416 \\
\addlinespace
\textsc{iii} & Institucionalidad concursal y auxiliares de la justicia
& Órganos del concurso y su fiscalización
& Estructura de la SUPERIR, categorías A y B de veedores y liquidadores, y control de cuentas \\
\addlinespace
\textsc{iv} & Procedimiento concursal de reorganización judicial ordinario
& Salvataje de la mediana y gran empresa
& Tramitación de la PFC ordinaria, quórums y financiamiento preferente \\
\addlinespace
\textsc{v} & Procedimiento concursal de reorganización simplificada de MIPEs
& Reestructuración ágil para la pequeña empresa
& Reemplazo de auditorías, veedores Categoría B y el reintento del artículo 286 S \\
\addlinespace
\textsc{vi} & Procedimiento administrativo de renegociación de la Persona Deudora
& Tramitación gratuita ante la SUPERIR
& Ingreso de boletas de honorarios, audiencias de determinación, renegociación y ejecución \\
\addlinespace
\textsc{vii} & Procedimiento concursal de liquidación ordinario
& Realización colectiva de activos empresariales
& Desasimiento, juicio de oposición, causales de liquidación forzosa y régimen recursivo restrictivo \\
\addlinespace
\textsc{viii} & Procedimiento concursal de liquidación simplificada
& Realización rápida para personas y MIPEs
& Entrega voluntaria, desformalización de la incautación, venta electrónica y fin de la cuenta tradicional \\
\addlinespace
\textsc{ix} & Régimen especial de acciones revocatorias concursales
& Ineficacia de actos lesivos a la masa
& Período sospechoso, acciones objetivas y subjetivas, y exclusión del abandono del procedimiento \\
\addlinespace
\textsc{x} & La paradoja del \emph{discharge} y la exclusión del Crédito con Aval del Estado
& Incompatibilidad de deudas de estudio con el perdón concursal
& Tensión entre universalidad y especialidad, y evolución de la jurisprudencia de la Corte Suprema \\
\addlinespace
\textsc{xi} & El incidente de mala fe del deudor concursal
& Límites éticos a la extinción de saldos insolutos
& Causales tasadas del artículo 169 A, sustanciación, sana crítica y recursos \\
\addlinespace
\textsc{xii} & Intersección del derecho concursal con el derecho laboral y penal económico
& Efectos en los contratos de trabajo y delitos societarios
& Artículo 163 bis del Código del Trabajo, finiquitos electrónicos y la Ley \Nro 21.595 \\
\addlinespace
\textsc{xiii} & Derecho concursal internacional e insolvencia transfronteriza
& Cooperación judicial y extraterritorialidad
& Ley Modelo UNCITRAL, determinación del COMI y análisis de los casos LATAM y Astaldi \\
\end{longtable}}

\cleardoublepage
